% =====================================================================
% Figure blocks for main.tex.
%
% PREAMBLE: add \usepackage{graphicx} -- main.tex does not currently load
% it, and neurips_2026.sty should not be relied on to.
%
% Upload the PDFs from docs/paper/figures/ to the Overleaf project.
% If you put them in a subfolder, add \graphicspath{{figures/}}.
%
%   fig1  ->  replaces the empty \fbox figure in Section 3.2
%   fig2  ->  new, in Section 3.1 or 5.1
%   fig3  ->  new, in Section 3.4
%   fig4  ->  replaces Tables tab:ens-gain, tab:ens-sig AND tab:asqa
%   figA1 ->  replaces Table tab:rank, moved to the appendix
%   figA2 ->  replaces Table tab:calib, moved to the appendix
%   figA3 ->  replaces Table tab:curve, moved to the appendix
%
% Regenerate with:
%   ./venv/bin/python -m open_ended_aggregation.analysis.pattern_cells
%   ./venv/bin/python -m open_ended_aggregation.analysis.budget_trace \
%       --qid 68__wikitables_composition__dev --json results/worked_example_hit.json
%   ./venv/bin/python -m open_ended_aggregation.analysis.budget_trace \
%       --qid 4__wikitables_composition__dev  --json results/worked_example_miss.json
%   ./venv/bin/python -m open_ended_aggregation.analysis.make_figures
% Each PDF ships a .manifest.json recording its source data and transforms.
% =====================================================================


% ---------------------------------------------------------------- fig 1 ---
\begin{figure}
  \centering
  \includegraphics[width=\linewidth]{fig1_worked_example.pdf}
  \caption{One held-out question, end to end, with every quantity of
  Section~\ref{sec:background} shown on real data. Rows are the pooled
  candidate set $C(q)$; a filled cell means that agent asserted the candidate,
  so the filled cells in a row are the support pattern $S(c)$ and the number
  inside is the within-output position $r_m(c)$; column headers give each
  agent's own output length $L_m(q)$; blank cells are the omissions that
  omission strength grades by $\log L_m(q)$. The scorer of
  Section~\ref{sec:slotA} orders the rows and the keep-count of
  Section~\ref{sec:slotB} draws the line. The takeaway is that support alone
  would not separate these rows cleanly, since \texttt{barchana} is asserted by
  four agents and \texttt{jajpur road} by one, yet the ordering that results
  keeps all seven gold answers and drops all seven non-answers. This question
  is favourable; Figure~\ref{fig:budget} shows one where the same rule misses.}
  \label{fig:worked}
\end{figure}


% ---------------------------------------------------------------- fig 2 ---
\begin{figure}
  \centering
  \includegraphics[width=\linewidth]{fig2_ceiling.pdf}
  \caption{The support-pattern partition on QAMPARI, which is what the ceiling
  of Section~\ref{sec:ceiling} is made of. Left: each marker is one support
  pattern $s$, positioned at its support count and its observed correctness
  rate, with area proportional to the claims it holds; the black line is the
  rate a rule reading only $|S(c)|$ can see. Any rule in the bounded class
  chooses one decision per marker, and multi-agent counting chooses one per
  column, so the vertical spread within a column is exactly the headroom
  reliability weighting has over counting. Right: where the claims actually
  are. The takeaway is that the spread is large where the claims are not:
  it reaches 59 points at $|S(c)|=4$, which holds 1,655 claims, and is 5.5
  points at $|S(c)|=1$, which holds 95,772. That is why a bound this loose in
  principle is worth only $+0.30$ F1 in practice.}
  \label{fig:ceiling}
\end{figure}


% ---------------------------------------------------------------- fig 3 ---
\begin{figure}
  \centering
  \includegraphics[width=\linewidth]{fig3_budget.pdf}
  \caption{How the keep-count is chosen, on two held-out questions. The solid
  curve is the plug-in expected F1 that Section~\ref{sec:slotB} maximises,
  computed from the predicted probabilities and $\hat{G}_q$ alone; the dashed
  curve is the true F1, which the rule never sees. Left: the argmax lands on
  the true optimum, and keeping every candidate instead would cost 33 F1.
  Right: the same rule cuts at 7 against a true optimum of 12 and ends below
  keeping everything. Note that $\hat{G}_q$ is wrong in both panels, by $+2.4$
  and $+1.1$ answers. The takeaway is that the rule needs the curve to have an
  interior maximum, not the gold-count estimate to be accurate, which is why
  substituting the true count is worth only $+0.55$ F1 overall.}
  \label{fig:budget}
\end{figure}


% ---------------------------------------------------------------- fig 4 ---
\begin{figure}
  \centering
  \includegraphics[width=\linewidth]{fig4_mirror.pdf}
  \caption{Gain over multi-agent counting for every agent subset of size 3 to
  6, with each point one ensemble and each bar the group median. Counting is
  the baseline and is drawn at zero. On QAMPARI the support-pattern ceiling
  sits above reliability weighting at every size and below the
  output-structure scorer at every size, so the scorer clears the bound rather
  than beating one weighting scheme; reliability weighting is significantly
  \emph{worse} than counting in 55 of 210 ensembles, while the scorer is above
  baseline in all 210. On ASQA the ordering reverses at every size. The
  takeaway is that the two benchmarks give mirror-image verdicts on the same
  two arms, and Table~\ref{tab:solo} predicts which verdict applies before
  either method is built.}
  \label{fig:mirror}
\end{figure}


% --------------------------------------------------------------- fig A1 ---
\begin{figure}
  \centering
  \includegraphics[width=\linewidth]{figA1_position.pdf}
  \caption{Within-output position as a correctness signal, restricted to solo
  assertions to remove the support confound and stratified by the asserting
  agent's own output length to remove the verbosity confound. Left: on lists of
  5 items or fewer the relationship runs the wrong way, and the usable gradient
  appears only on longer lists. Right: the same measurement as a rank
  correlation, with ASQA added; every ASQA response-length band falls inside
  $|\rho| < 0.05$, the range where QAMPARI shows no usable signal, and $79\%$
  of ASQA responses are 3 or 4 sentences. The takeaway is that the signal the
  method depends on is absent at ASQA's response lengths, so the ASQA loss in
  Figure~\ref{fig:mirror} is a property of the answer format rather than a
  failure of the template.}
  \label{fig:position}
\end{figure}


% --------------------------------------------------------------- fig A2 ---
\begin{figure}
  \centering
  \includegraphics[width=\linewidth]{figA2_calibration.pdf}
  \caption{Calibration of the cross-fitted claim scorer over 110,136 QAMPARI
  claims. Left: observed minus mean predicted correctness by
  predicted-probability bin. Right: how many claims fall in each bin.
  Calibration error is under one point across the five lowest bins, which hold
  $97\%$ of all claims, and grows monotonically to $-5.8$ points in the highest
  bin. The takeaway is that the error is concentrated in exactly the region the
  keep-count reads, so the budget overstates expected hits and keeps too much,
  and a globally fitted isotonic curve cannot fix it because such a curve is
  dominated by the bottom bin.}
  \label{fig:calibration}
\end{figure}


% --------------------------------------------------------------- fig A3 ---
\begin{figure}
  \centering
  \includegraphics[width=\linewidth]{figA3_learning_curve.pdf}
  \caption{Supervision cost of the claim scorer, on a fixed held-out test set
  with five independent training draws per size and two independent test
  splits. Left: the gain over counting is negative at 25 labelled questions and
  saturates by roughly 200 on both splits. Right: the range across the five
  draws collapses from $1.64$ to $0.15$ F1. The point at 527 is omitted from
  the right panel because it exhausts the training pool, leaving a single draw
  for which a range is undefined. The takeaway is that past roughly 200
  questions it no longer matters which questions were used for fitting, which
  is what one would expect if the coefficients describe how these models emit
  lists rather than the subject matter of any particular question.}
  \label{fig:curve}
\end{figure}
